\begingroup
\section{Node G0: compact global assembly and equality rigidity}
\label{module:G0}
\noindent\textit{Audited source: \nolinkurl{G0_global_wrapper.tex}.}
\subsection{Statement and normalization}

For $A>0$, let $\Pclass_N^{\rm conv}(A)$ be the class of bounded convex
planar polygons of area $A$ with exactly $N$ genuine sides, after
consecutive collinear sides are merged.  Let $R_N(A)$ be the regular
$N$-gon of area $A$.

\begin{theorem}[Global theorem conditional on J0]
\label{G0:thm:main}
Assume the active near-disk coercivity input in
\ref{G0:thm:J0}.  Then
For every $A>0$ there is $N_0$ such that, for all $N\ge N_0$ and all
$P\in\Pclass_N^{\rm conv}(A)$,
\begin{equation}\label{G0:eq:main}
 \lambda_1(P)\ge\lambda_1(R_N(A)).
\end{equation}
Equality holds if and only if $P$ is congruent to $R_N(A)$.
\end{theorem}

Put
\begin{equation}\label{G0:eq:Lambda}
 \LambdaF(\Omega)=\frac{|\Omega|}{\pi}\lambda_1(\Omega).
\end{equation}
It is invariant under dilations.  Hence it is enough to work at area
$\pi$; the general statement follows by the dilation
$P\mapsto(\pi/A)^{1/2}P$.

Let $\Kclass_{\le N}^{\rm conv}(\pi)$ be the Hausdorff-closed class of
area-$\pi$ convex polygons with at most $N$ active sides.  This class is
used only to take a global minimum.  The theorem itself concerns exactly
$N$ sides.

\subsection{The two retained local inputs}

We isolate the only two nonclassical inputs used in the global assembly.

\begin{theorem}[Strict transverse side activation; stage 132]
\label{G0:thm:activation}
Let $Q$ be a stationary convex $M$-gon for $\LambdaF$ in its active
stratum.  Insert a marked point in the relative interior of any side and
move it outward while keeping the two endpoints fixed.  For all
sufficiently small positive displacements this gives a convex
$(M+1)$-gon $Q_\varepsilon$ and
\begin{equation}\label{G0:eq:activation}
 \frac d{d\varepsilon}\LambdaF(Q_\varepsilon)
 \bigg|_{\varepsilon=0+}<0.
\end{equation}
The assertion is independent of all side lengths and mesh ratios.
\end{theorem}

The proof in stage 132 uses the stationary mass and barycenter identities
for the squared ground-state flux on a side, log-concavity of that flux,
and Hadamard's formula.  In particular, it is a first-variation theorem at
the lower-side point; it is not a mountain-pass deformation theorem.

\begin{obligation}[Active near-disk value coercivity; current J0 input]
\label{G0:thm:J0}
There are $\eta_0,S_0,c_0>0$ such that the following holds.  Let $P$ be an
area-$\pi$ active convex $N$-gon with positive normal fan
\begin{equation*}
 \alpha_i>0,\qquad \sum_i\alpha_i=2\pi,\qquad
 a=\frac{2\pi}{N},\qquad S=a+\max_i\alpha_i\le S_0.
\end{equation*}
Assume, after translation, that its radial graph is within $\eta_0$ of
the disk in $W^{1,\infty}$.  Let $h$ be its support vector in the fixed
translation complement of JP.  Put $h_0=h_\alpha\mathbf1$,
$A_0=A_\alpha(h_0)$, and
\(
s_\alpha(h)=A_\alpha'[h_0;h]/(2A_0).
\)
Normalize scale by
\[
 \widetilde h=\frac{h}{s_\alpha(h)},
 \qquad s_\alpha(\widetilde h)=1,
 \qquad \widetilde h=h_\alpha\mathbf1+z,
 \qquad z\in\mathsf X_\alpha.
\]
Shrinking the physical chart if necessary, $s_\alpha(h)>0$ and this
normalization stays in the same uniform near-disk chart.  Then
\begin{equation}\label{G0:eq:J0}
 \boxed{
 \lambda_1(P)-\lambda_1(R_N(\pi))
 \ge c_0\left\{
 \Jfour(\alpha)
 +\Gjp_\alpha(z-z_\alpha^\sharp)\right\},}
\end{equation}
where
\begin{equation}\label{G0:eq:J4}
 \Jfour(\alpha)=\sum_i\alpha_i^4-Na^4.
\end{equation}
The support metric is positive on the fixed affine scale section
$\mathsf X_\alpha$ after the fixed translation normalization.  Since
$\Phi_\alpha$ is scale invariant and the original $P$ has area $\pi$,
the JP/NR value inequality for $\widetilde h$ is exactly
\eqref{G0:eq:J0} for $\lambda_1(P)$.
Every constant is independent of the smallest physical, polar, and
conformal cells and of adjacent ratios.
\end{obligation}

For clarity, the conditional analytic chain behind input~\ref{G0:thm:J0} is
\begin{equation}\label{G0:eq:analytic-chain}
 \begin{gathered}
 \text{TR: exact physical trace, disk endpoint gap, value split}\\
 \downarrow\\
 \text{CT core}\longrightarrow\{\text{JP--EH},\text{JP--NS}\}\\
 \text{JP--NS+TR+CT core+QB+DF+DM}\longrightarrow\text{JP--DE}\\
 \text{JP--EH+JP--NS+JP--DE}\longrightarrow
 \{\text{JC--ALG},\text{JC--TL}\}\\
 \text{JC--ALG+JC--TL}\longrightarrow\text{JC--PL}
 \longrightarrow\text{JP--JC}\\
 \downarrow\\
 \text{JP--DE+JP--JC}\longrightarrow\text{NR}\longrightarrow\text{J0}.
 \end{gathered}
\end{equation}
No stationary Hessian is integrated on the nonstationary chart.

\subsection{Compact minimization in the closed side class}

\begin{lemma}[Compact spectral sublevels]\label{G0:lem:compact}
For fixed $N$ and $L<\infty$, the family
\begin{equation}\label{G0:eq:sublevel}
 \{P\in\Kclass_{\le N}^{\rm conv}(\pi):
       \lambda_1(P)\le L\}
\end{equation}
is Hausdorff precompact up to translations.  Its Hausdorff limits remain
in $\Kclass_{\le N}^{\rm conv}(\pi)$, and $\lambda_1$ is continuous along
such nondegenerate convex-domain convergence.  Consequently the infimum
of $\lambda_1$ over $\Kclass_{\le N}^{\rm conv}(\pi)$ is attained.
\end{lemma}

\begin{proof}
If $w(P)$ is the minimum width, containment in a strip and the
one-dimensional Poincar\'e inequality give
\[
 \lambda_1(P)\ge\frac{\pi^2}{w(P)^2}.
\]
Thus a spectral sublevel has a positive width bound.  For convex bodies,
fixed area and a positive minimum-width bound give a diameter bound.  To
see this without an implicit eccentricity lemma, let $[A,B]\subset P$ be
a diameter, $D=|A-B|$, and use it as the horizontal axis.  If $y_+$ and
$y_-$ are the extreme signed heights of $P$ above and below that axis,
the two triangles with base $[A,B]$ and apices at the corresponding
support points have disjoint interiors and lie in $P$.  Hence
\[
 |P|\ge \frac D2(y_+-y_-)
      \ge \frac D2 w(P),
 \qquad
 D\le \frac{2\pi}{w(P)}.
\]
After fixing, for example, the Steiner point at the origin, Blaschke
compactness then gives a Hausdorff convergent
subsequence.  Area is continuous for convex bodies, and having at most
$N$ genuine sides is closed (parametrize by at most $N$ cyclic vertices
and allow adjacent vertices to coalesce).  Dirichlet spaces of bounded convex
domains converge in the Mosco sense under this Hausdorff convergence, so
the first eigenvalue converges.  A minimizing sequence may be restricted
to the sublevel determined by any one competitor, proving attainment.
\end{proof}

\begin{proposition}[A closed-class minimizer loses no side]
\label{G0:prop:no-loss}
Every minimizer $Q_N$ of $\lambda_1$ over
$\Kclass_{\le N}^{\rm conv}(\pi)$ has exactly $N$ active sides.
\end{proposition}

\begin{proof}
After collinear sides are merged, suppose that $Q_N$ has $M<N$ genuine
sides.  The active $M$-gon stratum is a smooth finite-dimensional
manifold near $Q_N$ after the area and rigid gauges are fixed.  Since
$Q_N$ is a global minimizer, it is stationary on that stratum.
Theorem~\ref{G0:thm:activation} supplies a convex $(M+1)$-gon with strictly
smaller $\LambdaF$, hence with strictly smaller $\lambda_1$ at fixed
area.  It still belongs to $\Kclass_{\le N}^{\rm conv}(\pi)$, a
contradiction.
\end{proof}

\begin{remark}[Why SB disappears]
Stage 133 correctly observes that different one-side activation branches
need not be connected below the value of their common lower stratum.
Proposition~\ref{G0:prop:no-loss} does not connect them.  It selects a global
minimizer and needs only one strict feasible descent to contradict
minimality.  Therefore the side-branch bypass SB is not an open edge of
this proof.
\end{remark}

\subsection{Qualitative global-to-local entry}

Let $B$ be the area-$\pi$ disk and
$\lambda_B=\lambda_1(B)$.

\begin{lemma}[Regular polygons converge spectrally to the disk]
\label{G0:lem:regular-limit}
The area-$\pi$ regular polygons satisfy
\begin{equation}\label{G0:eq:regular-limit}
 R_N(\pi)\longrightarrow B\quad\text{in Hausdorff distance},
 \qquad
 \lambda_1(R_N(\pi))\longrightarrow\lambda_B.
\end{equation}
\end{lemma}

\begin{proof}
The inradius and circumradius of the area-normalized regular $N$-gon both
tend to one.  This proves Hausdorff convergence.  Spectral convergence is
the convex-domain stability used in Lemma~\ref{G0:lem:compact}.
\end{proof}

\begin{lemma}[Near-disk entry from a competing value]
\label{G0:lem:entry}
Let $N_j\to\infty$ and let
$P_j\in\Kclass_{\le N_j}^{\rm conv}(\pi)$ satisfy
\begin{equation}\label{G0:eq:competing}
 \lambda_1(P_j)\le\lambda_1(R_{N_j}(\pi)).
\end{equation}
Then, after translations,
\begin{equation}\label{G0:eq:support-limit}
 P_j\longrightarrow B\quad\text{in Hausdorff distance},\qquad
 \|h_{P_j}-1\|_{L^\infty(\mathbb S^1)}\longrightarrow0.
\end{equation}
If every $P_j$ is active and $\alpha_{\max}(P_j)$ denotes its largest
gap between consecutive outer normals, then
\begin{equation}\label{G0:eq:angle-limit}
 \alpha_{\max}(P_j)\longrightarrow0.
\end{equation}
Finally its radial function $r_j=1+\rho_j$ satisfies
\begin{equation}\label{G0:eq:radial-limit}
 \|\rho_j\|_{W^{1,\infty}(\mathbb S^1)}\longrightarrow0.
\end{equation}
\end{lemma}

\begin{proof}
Faber--Krahn and \eqref{G0:eq:competing}--\eqref{G0:eq:regular-limit} give
\[
 \lambda_B\le\lambda_1(P_j)
 \le\lambda_1(R_{N_j}(\pi))\longrightarrow\lambda_B.
\]
The sequence lies in one fixed spectral sublevel.  The proof of
Lemma~\ref{G0:lem:compact}, with no fixed side bound needed for the convex
compactness part, gives Hausdorff precompactness after translation.
Every limit has area $\pi$ and first eigenvalue $\lambda_B$.
Faber--Krahn rigidity makes the limit the disk.  Since every subsequential
limit is the disk, the whole sequence converges, and convergence of convex
support functions gives \eqref{G0:eq:support-limit}.

Suppose \eqref{G0:eq:angle-limit} failed.  Along a subsequence two adjacent
supporting normals would have a gap at least $\varepsilon>0$.  Their
support numbers tend uniformly to one by \eqref{G0:eq:support-limit}; their
intersection vertex in the bisector direction would therefore have
distance at least $\sec(\varepsilon/2)+o(1)>1$ from the origin.  This
contradicts Hausdorff convergence to $B$.

For a ray meeting the side with normal angle $\theta_i$ and support
number $h_i$,
\begin{equation}\label{G0:eq:radial-side}
 r_j(\theta)=\frac{h_i}{\cos(\theta-\theta_i)}.
\end{equation}
Hausdorff convergence and the inclusions
$(1-\varepsilon_j)B\subset P_j\subset(1+\varepsilon_j)B$ give
$r_j\to1$ uniformly, while support convergence gives $h_i\to1$
uniformly in $i$.  Thus, whenever the ray at $\theta$ meets side $i$,
\[
 \cos(\theta-\theta_i)=\frac{h_i}{r_j(\theta)}\longrightarrow1
\]
uniformly.  In particular $|\theta-\theta_i|\to0$; this also proves the
needed statement at one-sided sector endpoints.  Equation
\eqref{G0:eq:radial-side} and its one-sided derivative therefore give
\eqref{G0:eq:radial-limit}.
\end{proof}

\subsection{Contradiction and equality rigidity}

\begin{proof}[Proof of Theorem~\ref{G0:thm:main}]
By scale invariance, take $A=\pi$.

Assume first that the inequality fails for arbitrarily large side counts.
Then there are $N_j\to\infty$ and exact $N_j$-gons $P_j$ such that
\begin{equation}\label{G0:eq:strict-counter}
 \lambda_1(P_j)<\lambda_1(R_{N_j}(\pi)).
\end{equation}
Let $Q_j$ minimize $\lambda_1$ over
$\Kclass_{\le N_j}^{\rm conv}(\pi)$.  Lemma~\ref{G0:lem:compact} and
Proposition~\ref{G0:prop:no-loss} show that $Q_j$ exists and has exactly
$N_j$ active sides.  Moreover
\begin{equation}\label{G0:eq:min-counter}
 \lambda_1(Q_j)\le\lambda_1(P_j)
 <\lambda_1(R_{N_j}(\pi)).
\end{equation}
Lemma~\ref{G0:lem:entry} gives
\[
 \|\rho_{Q_j}\|_{W^{1,\infty}}\to0,\qquad
 S_j=\frac{2\pi}{N_j}+\alpha_{\max}(Q_j)\to0.
\]
Thus input~\ref{G0:thm:J0} applies for all sufficiently large $j$ and
gives
\[
 \lambda_1(Q_j)-\lambda_1(R_{N_j}(\pi))
 \ge c_0\{\Jfour(\alpha^{(j)})
 +\Gjp_{\alpha^{(j)}}(
 z_j-z_{\alpha^{(j)}}^\sharp)\}\ge0,
\]
contradicting \eqref{G0:eq:min-counter}.  Hence the inequality in
\eqref{G0:eq:main} holds for every sufficiently large $N$.

Suppose next that nonregular equality cases exist for arbitrarily large
$N$.  Choose $N_j\to\infty$ and nonregular exact $N_j$-gons $P_j$ with
\[
 \lambda_1(P_j)=\lambda_1(R_{N_j}(\pi)).
\]
The inequality just proved shows that this common value is the minimum
over the exact $N_j$-side class.  If the minimum over the closed class
with at most $N_j$ sides were smaller, its minimizer would have exactly
$N_j$ sides by Proposition~\ref{G0:prop:no-loss}, contradicting the proved
inequality.  Hence each $P_j$ is also a closed-class minimizer.
Lemma~\ref{G0:lem:entry} and input~\ref{G0:thm:J0} apply and force
\begin{equation}\label{G0:eq:defect-zero}
 \Jfour(\alpha^{(j)})=0,\qquad
 \Gjp_{\alpha^{(j)}}(
 z_j-z_{\alpha^{(j)}}^\sharp)=0.
\end{equation}
Strict convexity of $x^4$ under $\sum_i\alpha_i=2\pi$ makes the first
identity equivalent to $\alpha_i=2\pi/N_j$ for every $i$.  At the uniform
fan cyclic symmetry gives
$z_{\alpha^{(j)}}^\sharp=0$ on the exact quotient.  Positivity
of the support metric then makes the support vector constant modulo
translations and scale.  Thus $P_j$ is congruent to the regular
$N_j$-gon, a contradiction.

Both a strict failure and a nonregular equality can therefore occur only
for finitely many side counts.  Increasing $N_0$ proves the theorem.
\end{proof}
\endgroup
